\documentclass[12pt,final]{aunarAPA}

% ================== Paquetes básicos ==================
\usepackage[T1]{fontenc}
\usepackage[utf8]{inputenc}
\usepackage[spanish, es-tabla]{babel}

\usepackage{graphicx}
\usepackage{times}
\usepackage{subfiles}

\usepackage{amsmath,amssymb,amsfonts,amsthm,bm}
\usepackage{enumitem}
\usepackage{siunitx}
\usepackage{tocloft}
\usepackage{titlesec}
\usepackage{indentfirst}
\usepackage{makecell}
\usepackage[margin=2.2cm]{geometry}

% Bibliografía (APA con biber)
\usepackage[style=apa, backend=biber, language=spanish, maxcitenames=2, maxbibnames=99]{biblatex}
%\usepackage[style=apa]{biblatex}
\DeclareLanguageMapping{spanish}{spanish-apa}
\addbibresource{Biblio.bib}
%\usepackage[]{biblatex}


\usepackage{setspace}
\doublespacing
\usepackage[document]{ragged2e}
\let\cite\parencite

% ================== Tablas ==================
\usepackage{tabularx}
\usepackage{booktabs}
\usepackage{longtable}
\usepackage{multirow}
\usepackage{array}
\usepackage{caption}
\usepackage{threeparttablex}

% Ajustes globales de tablas
\renewcommand{\arraystretch}{1.22}
\setlength{\tabcolsep}{4pt}

% ---- Formato de títulos de tablas ----
\captionsetup[table]{%
  name=Tabla,
  labelfont=bf,
  textfont=it,
  singlelinecheck=false,
  justification=raggedright
}
\captionsetup[longtable]{name=Tabla}
\makeatletter
\renewcommand{\fnum@table}{\tablename~\thetable}
\makeatother

% ================== Índices ==================
\renewcommand{\contentsname}{Tabla de contenido}
\renewcommand{\cfttoctitlefont}{\hfill\bfseries\normalsize}
\renewcommand{\cftaftertoctitle}{\hfill}

\renewcommand{\cftchapfont}{\bfseries}
\renewcommand{\cftchappagefont}{}
\renewcommand{\cftsecfont}{}
\renewcommand{\cftsecpagefont}{}
\renewcommand{\cftsubsecfont}{}
\renewcommand{\cftsubsecpagefont}{}
\renewcommand{\cftsubsubsecfont}{}
\renewcommand{\cftsubsubsecpagefont}{}

\cftsetindents{chapter}{0pt}{3.5em}
\cftsetindents{section}{1.8em}{3.2em}
\cftsetindents{subsection}{3.8em}{3.2em}
\cftsetindents{subsubsection}{5.5em}{3.2em}

% ================== Lista de Figuras y Tablas ==================
\renewcommand{\listfigurename}{Lista de figuras}
\renewcommand{\listtablename}{Lista de tablas}
\renewcommand{\cftloftitlefont}{\hfill\bfseries\normalsize}
\renewcommand{\cftafterloftitle}{\hfill}
\renewcommand{\cftlottitlefont}{\hfill\bfseries\normalsize}
\renewcommand{\cftafterlottitle}{\hfill}
\renewcommand{\cftfigfont}{}
\renewcommand{\cftfigpagefont}{}
\renewcommand{\cfttabfont}{}
\renewcommand{\cfttabpagefont}{}

\sisetup{locale=DE, table-number-alignment=center}

% ================== Títulos (mismo tamaño, negrita) ==================
\titleformat{\chapter}[hang]{\bfseries\normalsize}{\thechapter.}{0.5em}{}
\titlespacing*{\chapter}{0pt}{0.75ex}{0.75ex}

\titleformat{\section}[hang]{\bfseries\normalsize}{\thesection}{0.5em}{}
\titlespacing*{\section}{0pt}{0.75ex}{0.5ex}

\titleformat{\subsection}[hang]{\bfseries\normalsize}{\thesubsection}{0.5em}{}
\titlespacing*{\subsection}{1.5em}{0.75ex}{0.5ex}

\titleformat{\subsubsection}[hang]{\bfseries\normalsize}{\thesubsubsection}{0.5em}{}
\titlespacing*{\subsubsection}{3em}{0.75ex}{0.5ex}

% ================== “Capítulos pequeños” ==================
\newcommand{\smallchapter}[1]{%
  \clearpage
  \vspace*{1cm}
  \begin{center}\textbf{\normalsize #1}\end{center}
  \vspace{0.5cm}
  \addcontentsline{toc}{chapter}{#1}
}

% ================== Hipervínculos ==================
\usepackage[hidelinks]{hyperref}

% ================== Metadatos aunarAPA ==================
\institucion{Universidad Autónoma de Bucaramanga}
\facultad{Ingeniería}
\programa{Ingeniería Biomédica}
\ciudad{Bucaramanga}
\yeardef{2025}

% ================== Documento ==================
\begin{document}

% ----------- PORTADA ÚNICA -----------
\begin{titlepage}
\centering
\vspace*{1cm}
{\bfseries\normalsize
Apósito con Geometría Auxética para la Epitelización de Quemaduras de Segundo Grado por Fricción\par}

\vspace{3.5cm}
{\bfseries\normalsize Autores:}\\[0.6cm]
{\normalsize
Nombre1\\
Nombre2\\
Nombre3}

\vspace{2.2cm}
{\bfseries\normalsize Director:} PhD. Mateo Escobar Jaramillo

\vspace{3.5cm}
{\normalsize
Universidad Autónoma de Bucaramanga\\
Facultad de Ingeniería\\
Programa de Ingeniería Biomédica}

\vfill
{\normalsize 2025}
\end{titlepage}
% ----------- FIN PORTADA -----------

% ================== Configuración general ==================
\setcounter{secnumdepth}{4}
\setcounter{tocdepth}{3}
\setlength{\parindent}{1.27cm}

% ================== Preliminares e índices ==================
\clearpage
\pagenumbering{Roman}
\setcounter{page}{1} % Portada = I

\newpage
\tableofcontents
\newpage
\listoffigures
\newpage
\listoftables

\newpage 
\subfile{./1agradecimientos/agradecimientos}   % II
\newpage 
\subfile{./2dedicatoria/dedicatoria}       % III
\newpage 
\subfile{./3resumen/resumen}           % IV
\newpage 
\subfile{./4abstract/abstract}          % V

\clearpage
\pagenumbering{arabic} % volver a numeración normal

% ================== Capítulos ==================
\newpage 
\subfile{./5cap1/capitulo1}

\newpage 
\subfile{./6cap2/capitulo2}

\newpage 
\subfile{./7cap3/capitulo3}

\newpage 
\subfile{./8cap4/capitulo4}

\newpage 
\subfile{./9cap5/capitulo5}

% ================== Cierres ==================
\newpage 
\subfile{./10conclusiones/conclusiones}
\newpage 
\subfile{./11recomendaciones/recomendaciones}
\newpage 
\subfile{./12referencias/referencias}
\newpage 
\subfile{./13anexos/anexos}

\end{document}

